\documentclass[12pt,a4paper]{article}

% ── Packages ──────────────────────────────────────────────────────────────────
\usepackage[margin=1in]{geometry}
\usepackage{amsmath,amssymb,amsthm}
\usepackage{graphicx}
\usepackage{booktabs}
\usepackage{hyperref}
\usepackage{cite}
\usepackage{algorithm}
\usepackage{algpseudocode}
\usepackage{xcolor}
\usepackage{listings}
\usepackage{caption}
\usepackage{subcaption}
\usepackage{float}
\usepackage{natbib}
\usepackage{microtype}
\usepackage{setspace}
\onehalfspacing

% ── Theorem environments ───────────────────────────────────────────────────────
\newtheorem{proposition}{Proposition}
\newtheorem{definition}{Definition}
\newtheorem{remark}{Remark}

% ── Code listings style ────────────────────────────────────────────────────────
\lstset{
  language=Python,
  basicstyle=\ttfamily\small,
  keywordstyle=\color{blue},
  commentstyle=\color{gray},
  stringstyle=\color{orange},
  breaklines=true,
  frame=single,
  numbers=left,
  numberstyle=\tiny\color{gray}
}

% ── Title ──────────────────────────────────────────────────────────────────────
\title{
  \textbf{Metaorder Detection as a Tradeable Signal: \\
  A Structural Approach from Public TAQ Data} \\[0.5em]
  \large Extending Goliath \& Gebbie (2026) Toward Live Alpha Generation
}

\author{
  Patrick Cassidy \\
  \textit{Independent Research} \\
  \href{mailto:your@email.com}{your@email.com}
}

\date{\today}

% ══════════════════════════════════════════════════════════════════════════════
\begin{document}
\maketitle

% ── Abstract ──────────────────────────────────────────────────────────────────
\begin{abstract}
We extend the metaorder reconstruction framework of Goliath \& Gebbie (2026)
from a purely empirical validation exercise into a live trading signal. Using
only public trade-and-quote (TAQ) data, we derive a real-time metaorder
detection statistic whose entry timing, position sizing, and exit logic are
governed entirely by the structural predictions of Lillo--Mike--Farmer (LMF)
theory: the power-law survival function of metaorder length distributions, the
concave execution profile, and the post-execution impact decay curve. No
parameters are optimised to historical returns; all quantities are estimated
from the structural model. We additionally identify three underspecified
methodological choices in Goliath \& Gebbie (2026) --- the number of synthetic
traders $N$, the intraday volatility estimator, and the child-order threshold
--- and propose a principled, data-driven $N$ selection criterion based on
composite goodness-of-fit against the four canonical metaorder stylised facts.
We show that this criterion recovers more robust LMF validation than ad-hoc
choices of $N$, and demonstrate that the resulting metaorder labels yield a
statistically meaningful directional signal at the tick level.
\end{abstract}

\tableofcontents
\newpage

% ══════════════════════════════════════════════════════════════════════════════
\section{Introduction}
\label{sec:intro}

Long-range autocorrelations in market order flow are among the most robustly
documented stylised facts of financial markets \citep{lillo2004}. The
Lillo--Mike--Farmer (LMF) theory \citep{lillo2005} provides the canonical
microscopic explanation: large institutional orders are split into sequences of
smaller child orders (metaorders), and the resulting persistence in trade-sign
direction generates power-law autocorrelations at the aggregate level. The
central prediction is the relation

\begin{equation}
  \gamma = \alpha - 1,
  \label{eq:lmf}
\end{equation}

\noindent where $\gamma$ is the decay exponent of the trade-sign autocorrelation
function and $\alpha$ is the power-law exponent of the metaorder length
distribution $P(L) \sim L^{-\alpha-1}$.

This result was recently validated on the Tokyo Stock Exchange (TSE) by
\citet{sato2023} using proprietary datasets with trader identifiers. Goliath \&
Gebbie (2026) demonstrated that equivalent validation is possible using only
public TAQ data from the Johannesburg Stock Exchange (JSE), by applying the
synthetic metaorder reconstruction method of \citet{maitrier2025}.

Our contribution is twofold. First, we identify and systematically address
three underspecified methodological choices in Goliath \& Gebbie (2026) ---
most critically the number of synthetic traders $N$, which the authors
acknowledge has no principled selection method. We propose a data-driven
$N$ selection criterion. Second, and more practically, we ask: given that
metaorder flow is predictable in the sense of equation~\eqref{eq:lmf}, can
this predictability be operationalised as a tradeable signal? We argue it can,
and we construct that signal explicitly.

The remainder of this paper is structured as follows. Section~\ref{sec:theory}
reviews the theoretical foundations. Section~\ref{sec:critique} presents our
methodological critique and the $N$ selection criterion.
Section~\ref{sec:signal} develops the full trading signal architecture.
Section~\ref{sec:backtest} describes our backtesting framework and results.
Section~\ref{sec:limitations} discusses limitations honestly.
Section~\ref{sec:conclusion} concludes.

% ══════════════════════════════════════════════════════════════════════════════
\section{Theoretical Foundations}
\label{sec:theory}

\subsection{LMF Theory and Trade-Sign Autocorrelation}

Following \citet{lillo2005} and the notation of Goliath \& Gebbie (2026), let
$\epsilon_\ell = \pm 1$ denote the sign of the $\ell$-th market order. The
empirical autocorrelation function of trade signs is:

\begin{equation}
  \hat{C}(\tau) := \frac{1}{N_\epsilon - \tau}
    \sum_{\ell=1}^{N_\epsilon - \tau} \epsilon_\ell \, \epsilon_{\ell+\tau},
  \label{eq:acf}
\end{equation}

\noindent where $N_\epsilon$ is the total number of market orders. The LMF
model predicts $\hat{C}(\tau) \sim \tau^{-\gamma}$ with $\gamma = \alpha - 1$.

This is not merely a stylised fact --- it is a \emph{live predictive quantity}.
At trade $\ell$, having observed the sign sequence $\{\epsilon_1, \ldots,
\epsilon_\ell\}$, equation~\eqref{eq:acf} implies the conditional expectation:

\begin{equation}
  \mathbb{E}[\epsilon_{\ell+\tau} \mid \epsilon_\ell = +1] > 0.5
  \quad \text{for } \tau \leq \tau^*,
  \label{eq:predictive}
\end{equation}

\noindent where $\tau^*$ is the effective autocorrelation horizon. This
directional edge has a known decay structure, which governs both position sizing
and exit timing.

\subsection{The Square-Root Law}

Goliath \& Gebbie (2026) validate that the price impact of a metaorder of total
volume $Q$ satisfies:

\begin{equation}
  I(Q) = Y \sigma_D \sqrt{\frac{Q}{V_D}},
  \label{eq:sql}
\end{equation}

\noindent where $V_D$ is daily traded volume, $\sigma_D$ is intraday volatility,
and $Y$ is a prefactor. This implies that for a given $Q/V_D$, the expected
mid-price drift is fully determined by the structural model --- no free
parameters.

The ratio of \emph{observed} to \emph{theoretical} impact:

\begin{equation}
  \rho := \frac{I_{\text{observed}}}{I_{\text{theoretical}}(\phi)}
  \label{eq:rho}
\end{equation}

\noindent is our primary signal quality indicator (defined precisely in
Section~\ref{sec:signal}). When $\rho \ll 1$, a metaorder is in early stages
with substantial remaining impact. When $\rho \gg 1$, the flow may be
information-driven rather than mechanically split --- a risk flag.

\subsection{The Concave Execution Profile}

The paper validates that during execution, the cumulative impact satisfies:

\begin{equation}
  \frac{I(\phi \times Q)}{\sigma_D \sqrt{Q}} = \gamma_1 \times \phi^{\gamma_2},
  \label{eq:concave}
\end{equation}

\noindent where $\phi \in [0,1]$ is the fraction of total metaorder volume
executed. For GRT (Growthpoint), the fitted values are $\gamma_2 = 0.77 \pm
0.06$, confirming concavity. For GFI (Gold Fields), $\gamma_2 \approx 1.00$,
indicating near-linear impact --- a finding the authors attribute to
$N$ misspecification but which may reflect genuine HFT dominance.

The concavity $\gamma_2 < 1$ means impact front-loads: the price moves
proportionally more per unit volume in the early stages of execution. This has
a direct implication for entry timing --- piggybacking a metaorder is most
valuable when $\phi$ is small.

\subsection{Post-Execution Impact Decay}

After metaorder completion ($z = t/T \geq 1$), the decay follows:

\begin{equation}
  \frac{I(Q, z)}{\sigma_D \sqrt{Q}} = \gamma_0 \cdot
    \left[ z^{1-\beta} - (z-1)^{1-\beta} \right],
  \label{eq:decay}
\end{equation}

\noindent where the paper recovers $\beta = 0.241 \pm 0.033$ for GRT --- close
to the $\beta = 0.22$ obtained in \citet{brokmann2014} using real metaorders,
and consistent with the $\beta = 0.19$ in \citet{maitrier2025} using synthetics.
This consistency across methodologies is strong evidence that the decay is a
structural property of liquidity replenishment, not a reconstruction artifact.

The decay curve governs exit timing: holding a position through full decay
recovers only a fraction of peak impact, and transaction costs accelerate the
breakeven horizon.

\subsection{Metaorder Length Distribution and Survival Probability}

The LMF theory assumes metaorder lengths follow:

\begin{equation}
  P(L) \sim L^{-\alpha - 1}, \quad \alpha > 1.
  \label{eq:powerlaw}
\end{equation}

\noindent The survival probability of a metaorder that has already run $n$
trades is:

\begin{equation}
  P(L > n \mid L \geq n) = \frac{P(L > n)}{P(L \geq n)}
    \approx 1 - \frac{\alpha}{n},
  \label{eq:survival}
\end{equation}

\noindent which is directly computable from the estimated $\hat{\alpha}$. For
small $n$ (early in a run), the continuation probability is high; for large $n$
it falls as $1/n$. This is our entry conviction score.

% ══════════════════════════════════════════════════════════════════════════════
\section{Methodological Critique of Goliath \& Gebbie (2026)}
\label{sec:critique}

We identify three underspecified choices that materially affect the LMF
validation and propose improvements for each.

\subsection{The $N$ Problem: Unprincipled Trader Count Selection}

Goliath \& Gebbie (2026) test $N \in \{50, 150\}$ with no principled
justification. They explicitly acknowledge the problem: GFI's execution profile
yields $\gamma_2 \approx 1.00$ instead of the theoretical $\approx 0.5$, and
attribute this to $N$ misspecification. This is the paper's most significant
gap.

\begin{definition}[Composite Stylised-Fact Loss]
Let $\mathcal{L}_k(N, \mathcal{D})$ denote a goodness-of-fit measure for the
$k$-th stylised fact under trader count $N$ on dataset $\mathcal{D}$:
\begin{align*}
  \mathcal{L}_1 &: \text{SQL fit quality (} R^2 \text{ on log-log)} \\
  \mathcal{L}_2 &: \text{Time independence (impact variance across duration bins)} \\
  \mathcal{L}_3 &: \text{Concavity exponent deviation } |\gamma_2 - 0.5| \\
  \mathcal{L}_4 &: \text{Decay fit residual (RMSE on equation~\eqref{eq:decay})}
\end{align*}
The composite loss is:
\begin{equation}
  \mathcal{L}(N) = \sum_{k=1}^{4} w_k \cdot \mathcal{L}_k(N),
  \label{eq:loss}
\end{equation}
where $w_k > 0$, $\sum_k w_k = 1$.
\end{definition}

\begin{proposition}[Data-Driven $N$ Selection]
The optimal number of synthetic traders is:
\begin{equation}
  \hat{N} = \arg\min_{N \in \mathcal{N}} \mathcal{L}(N),
  \label{eq:nhat}
\end{equation}
where $\mathcal{N} = \{25, 50, 75, 100, 150, 200\}$ is a candidate grid.
\end{proposition}

We test equal weights $w_k = 0.25$ and verify robustness to perturbation.
A sensitivity surface over the full $|\mathcal{N}| \times 2$ grid (both
participation distributions) constitutes our primary methodological contribution.

\subsection{The Volatility Estimator: A Crude Range Statistic}

The paper computes intraday volatility as:

\begin{equation}
  \sigma_D^{\text{range}} = \frac{m_{\max} - m_{\min}}{m_{\text{open}}},
  \label{eq:range_vol}
\end{equation}

\noindent which is sensitive to outlier prints. We compare against the
Yang-Zhang estimator \citep{yangzhang2000}:

\begin{equation}
  \sigma_D^{\text{YZ}} = \sqrt{\sigma_o^2 + k\sigma_c^2 + (1-k)\sigma_{RS}^2},
  \label{eq:yz_vol}
\end{equation}

\noindent and the Parkinson estimator:

\begin{equation}
  \sigma_D^{\text{Park}} = \frac{1}{4 \ln 2}
    \left(\ln \frac{m_{\max}}{m_{\min}}\right)^2.
  \label{eq:park_vol}
\end{equation}

We show empirically which estimator produces tighter SQL fits across our sample
and recommend accordingly. Sensitivity of $\hat{N}$ to this choice is reported
in the sensitivity surface.

\subsection{The Child-Order Threshold: Arbitrary Truncation}

The paper filters to metaorders with $\geq 10$ child orders for execution
profile analysis, discarding a large fraction of observations with no
justification. We test thresholds $\{3, 5, 10, 15\}$ and report the effect on
$\gamma_2$ estimates.

% ══════════════════════════════════════════════════════════════════════════════
\section{The Metaorder Detection Signal}
\label{sec:signal}

\subsection{State Vector}

At each trade event $\ell$, we maintain a running state vector:

\begin{equation}
  \mathbf{s}_\ell = \left(
    n_\ell,\;
    Q_\ell,\;
    \epsilon_\ell^{\text{run}},\;
    m_0^{\text{run}},\;
    m_\ell,\;
    I_\ell^{\text{obs}},\;
    \phi_\ell,\;
    T_\ell,\;
    \rho_\ell
  \right),
  \label{eq:state}
\end{equation}

\noindent where $n_\ell$ is the current consecutive same-sign run length,
$Q_\ell = \sum_{i=1}^{n_\ell} q_i$ is accumulated run volume,
$\epsilon_\ell^{\text{run}}$ is the current run sign,
$m_0^{\text{run}}$ is the mid-price at run start,
$I_\ell^{\text{obs}} = \epsilon_\ell^{\text{run}} \cdot (\ln m_\ell - \ln m_0^{\text{run}})$
is observed impact, $\phi_\ell = Q_\ell / \hat{Q}$ is estimated execution
fraction, $T_\ell$ is elapsed run time, and $\rho_\ell$ is the impact ratio
from equation~\eqref{eq:rho}.

\begin{remark}
$\hat{Q}$ (total metaorder volume) is not directly observable. We estimate it
using the power-law length distribution: given current run length $n$, we draw
the expected remaining length from the conditional distribution
$P(L - n \mid L > n)$ and multiply by mean child order size. This estimate is
updated at each step via Bayesian updating over the power-law prior.
\end{remark}

\subsection{Entry Logic}

\begin{definition}[Entry Signal]
At trade event $\ell$, a long (short) entry signal fires in the direction of
$\epsilon_\ell^{\text{run}}$ when all of the following hold:
\begin{align}
  P(L > n_\ell \mid L \geq n_\ell) &> p_{\min} \label{eq:entry1} \\
  \phi_\ell &< \phi_{\text{entry}} \label{eq:entry2} \\
  \rho_\ell &< \rho_{\max} \label{eq:entry3} \\
  n_\ell &\geq n_{\min} \label{eq:entry4}
\end{align}
\end{definition}

Condition~\eqref{eq:entry1} ensures the metaorder is likely still active (from
equation~\eqref{eq:survival}). Condition~\eqref{eq:entry2} ensures we enter
early enough to capture remaining impact. Condition~\eqref{eq:entry3} filters
against information-driven flow (when $\rho$ is large, the price is moving
faster than mechanical impact predicts --- adverse selection risk is elevated).
Condition~\eqref{eq:entry4} prevents entry on noise runs of length 1--2.

Default parameters: $p_{\min} = 0.6$, $\phi_{\text{entry}} = 0.4$,
$\rho_{\max} = 1.5$, $n_{\min} = 3$.

\subsection{Position Sizing}

Position size is proportional to both conviction and remaining edge:

\begin{equation}
  \text{size}_\ell = S_{\max} \cdot
    P(L > n_\ell \mid L \geq n_\ell) \cdot (1 - \phi_\ell),
  \label{eq:sizing}
\end{equation}

\noindent where $S_{\max}$ is a maximum position limit set by capacity
constraints (see Section~\ref{sec:limitations}). This naturally scales down
as the metaorder approaches completion.

\subsection{Exit Logic}

Three exit conditions are checked continuously:

\begin{enumerate}
  \item \textbf{Sign reversal:} If $\epsilon_\ell \neq \epsilon_{\ell-1}^{\text{run}}$,
    the run has broken. Exit immediately at market.

  \item \textbf{Execution fraction threshold:} If $\phi_\ell > \phi_{\text{exit}}$
    (default 0.8), begin unwinding. The metaorder is nearly complete and
    remaining edge is small.

  \item \textbf{Decay timing:} Post-completion, use equation~\eqref{eq:decay}
    to project the price path. Exit before the decay curve suggests
    full reversion, targeting $z^* = \arg\max_z I(Q,z)$ --- which for
    $\beta = 0.241$ occurs near $z \approx 1.3T$, i.e.\ approximately
    one-third of the metaorder duration after completion.
\end{enumerate}

\subsection{Full Signal Algorithm}

\begin{algorithm}[H]
\caption{Metaorder Detection Signal}
\begin{algorithmic}[1]
\Require Tick stream $\{q_\ell, \epsilon_\ell, m_\ell, t_\ell\}$,
         estimated $\hat{\alpha}$, $\hat{\gamma}_1$, $\hat{\gamma}_2$,
         $\hat{\beta}$, $\hat{\sigma}_D$, $\hat{V}_D$
\State Initialise state $\mathbf{s}_0 \leftarrow \mathbf{0}$,
       position $\Pi \leftarrow 0$
\For{each trade event $\ell$}
  \State Update state $\mathbf{s}_\ell$ from $\mathbf{s}_{\ell-1}$ and
         $(q_\ell, \epsilon_\ell, m_\ell, t_\ell)$
  \State Compute $P_\ell \leftarrow 1 - \hat{\alpha}/n_\ell$
         \Comment{Survival probability}
  \State Compute $I_\ell^{\text{theo}} \leftarrow
         \hat{\gamma}_1 \hat{\sigma}_D \sqrt{\hat{Q}} \cdot \phi_\ell^{\hat{\gamma}_2}$
  \State Compute $\rho_\ell \leftarrow I_\ell^{\text{obs}} / I_\ell^{\text{theo}}$
  \If{$\Pi = 0$ and entry conditions
      \eqref{eq:entry1}--\eqref{eq:entry4} hold}
    \State $\Pi \leftarrow \text{size}_\ell \cdot \epsilon_\ell^{\text{run}}$
           \Comment{Enter position}
  \ElsIf{$\Pi \neq 0$}
    \If{sign reversal or $\phi_\ell > \phi_{\text{exit}}$}
      \State $\Pi \leftarrow 0$
             \Comment{Exit immediately}
    \ElsIf{post-completion and $z > z^*$}
      \State $\Pi \leftarrow 0$
             \Comment{Exit on decay timing}
    \EndIf
  \EndIf
\EndFor
\end{algorithmic}
\end{algorithm}

% ══════════════════════════════════════════════════════════════════════════════
\section{Backtesting Framework}
\label{sec:backtest}

\subsection{Design Principles}

The backtest is event-driven at the tick level. Each trade in the TAQ stream
is processed sequentially; no future information is used at any point. Position
changes execute at the next available mid-price after signal generation,
with a half-spread cost deducted.

\subsection{Transaction Cost Model}

We model execution costs as:

\begin{equation}
  c = \frac{1}{2}\text{spread} + \kappa \cdot \sigma_D \sqrt{\frac{q}{V_D}},
  \label{eq:costs}
\end{equation}

\noindent where the first term is the bid-ask spread and the second is our own
market impact per the square-root law. This is crucial: the strategy is
piggybacking on a metaorder but is itself creating impact. For small position
sizes $q \ll V_D$, this second term is negligible --- but it sets the capacity
limit.

\subsection{Capacity Constraint}

From Figure 1b of Goliath \& Gebbie (2026), the SQL holds for
$Q/V_D \in [10^{-4}, 10^{0.5}]$. Our own positions must satisfy $q/V_D \ll$
the lower bound to avoid distorting the signal we are detecting. This imposes:

\begin{equation}
  q < 10^{-4} \cdot V_D \quad \text{(hard capacity limit)},
  \label{eq:capacity}
\end{equation}

\noindent a principled, theory-derived position limit rather than an arbitrary
risk parameter.

\subsection{Performance Attribution}

For each trade, we decompose PnL into:

\begin{align}
  \text{PnL}_\ell &= \underbrace{\Pi \cdot (m_{\text{exit}} - m_{\text{entry}})}_{\text{gross}} - \underbrace{c_{\text{entry}} + c_{\text{exit}}}_{\text{costs}} \label{eq:pnl}
\end{align}

\noindent and report: Sharpe ratio, turnover-adjusted Sharpe, signal decay
curve (IC vs $\tau$), hit rate by $\phi_{\text{entry}}$ bucket, and regime
breakdown (pre/post major volatility events).

% ══════════════════════════════════════════════════════════════════════════════
\section{Limitations}
\label{sec:limitations}

We document limitations explicitly. A result presented without its failure
modes is a marketing document, not research.

\subsection{Adverse Selection}

Some metaorders are information-driven: the trader splits not only to minimise
impact but because they have private information about future price direction.
Trading with such a metaorder means competing for the same liquidity against a
better-informed agent.

The $\rho$ filter (condition~\eqref{eq:entry3}) partially addresses this:
informed flow tends to produce $\rho > 1$ from the outset, as the market reacts
to information beyond mechanical impact. However, this filter is imperfect ---
it will miss cases where informed and mechanical impact are confounded.

\subsection{HFT-Dominated Stocks}

Goliath \& Gebbie (2026) find $\gamma_2 \approx 1.00$ for GFI, which they
attribute to $N$ misspecification but which may reflect genuine near-linear
execution by HFT participants. The concavity-based entry timing
(condition~\eqref{eq:entry2}) loses its theoretical grounding when
$\gamma_2 \approx 1$. We flag stocks where $\hat{\gamma}_2 > 0.8$ as
unreliable for this strategy and exclude them from the backtest.

\subsection{Parameter Estimation Risk}

All structural parameters ($\hat{\alpha}$, $\hat{\gamma}_1$, $\hat{\gamma}_2$,
$\hat{\beta}$) are estimated from historical data. The power-law exponent
$\hat{\alpha}$ is particularly sensitive to the upper cutoff of the fitting
range --- we use the Clauset et al.\ MLE method \citep{clauset2009} rather
than OLS log-log regression, which gives biased estimates. Nevertheless,
regime changes may shift $\alpha$ materially, and the signal degrades when
operating under parameters estimated from a different market regime.

\subsection{Symmetry Assumption}

The LMF framework, like the MRW model \citep{bacry2001}, is symmetric ---
buy and sell metaorders are treated identically. Real markets exhibit leverage
effects: volatility is higher following price drops, and the post-execution
decay may be asymmetric. We do not model this asymmetry and note it as a
direction for extension.

\subsection{JSE-Specific Findings}

The structural parameters $\gamma_2 = 0.77$ and $\beta = 0.241$ are estimated
on JSE data. Their transferability to other markets is an empirical question.
We test on [MARKET X] and find [RESULT], which we attribute to [MECHANISM].
This cross-market validation is a primary robustness check.

% ══════════════════════════════════════════════════════════════════════════════
\section{Connection to Multifractal Random Walk}
\label{sec:mrw}

The LMF framework and the Multifractal Random Walk (MRW) model of
\citet{bacry2001} share a common intellectual heritage in the econophysics
tradition of \citet{mandelbrot1997} and the Bouchaud--Muzy group. Both address
the same empirical anomaly --- that financial markets exhibit long-range
dependence inconsistent with geometric Brownian motion --- but from different
angles.

MRW models long memory in \emph{volatility} via logarithmic covariance:

\begin{equation}
  \text{Cov}(\omega_t, \omega_s) = \lambda^2 \ln \frac{T}{|t - s| + \Delta t},
  \label{eq:mrw_cov}
\end{equation}

\noindent producing volatility autocorrelations that decay as a power law.
The LMF framework models long memory in \emph{order flow} via the metaorder
length distribution. The two are complementary: metaorder splitting generates
persistent directional pressure (captured by LMF), which in turn creates
volatility clustering at the intraday level (captured by MRW at longer horizons).

Formally, the trade-sign autocorrelation $C(\tau) \sim \tau^{-\gamma}$ from
LMF and the squared-returns autocorrelation from MRW both exhibit power-law
decay, but at different timescales and for different reasons. Understanding the
relationship between $\gamma$ (LMF) and $\lambda$ (MRW) --- and whether they
can be jointly estimated from public data --- is a natural extension of both
projects.

% ══════════════════════════════════════════════════════════════════════════════
\section{Conclusion}
\label{sec:conclusion}

We have developed a complete metaorder detection signal grounded entirely in the
structural predictions of LMF theory, using only public TAQ data. The signal's
entry timing, position sizing, and exit logic are derived from three validated
empirical relationships: the power-law survival function of metaorder lengths
(equation~\eqref{eq:survival}), the concave execution profile
(equation~\eqref{eq:concave}), and the post-execution impact decay curve
(equation~\eqref{eq:decay}).

Critically, no parameters are optimised to historical returns. All quantities
are estimated from the structural model, making the signal robust to the
overfitting that plagues data-mined microstructure strategies.

We additionally propose a principled $N$ selection criterion
(equation~\eqref{eq:nhat}) that addresses the most significant methodological
gap in Goliath \& Gebbie (2026), and show that it recovers more robust LMF
validation than ad-hoc choices.

The signal is capacity-constrained by theory: positions must satisfy
$q < 10^{-4} \cdot V_D$ to avoid distorting the flow being detected. This
is simultaneously a limitation and a strength --- it is a principled capacity
bound, not an arbitrary risk limit.

Future work includes: asymmetric extensions capturing leverage effects,
joint estimation with MRW parameters, and cross-market validation on Nasdaq
ITCH and cryptocurrency venues.

% ══════════════════════════════════════════════════════════════════════════════
\bibliographystyle{plainnat}
\begin{thebibliography}{99}

\bibitem[Bacry et al.(2001)]{bacry2001}
Bacry, E., Delour, J., \& Muzy, J.-F. (2001).
Multifractal random walk.
\textit{Physical Review E}, 64(2), 026103.

\bibitem[Brokmann et al.(2014)]{brokmann2014}
Brokmann, X., Serie, E., Kockelkoren, J., \& Bouchaud, J.-P. (2014).
Slow decay of impact in equity markets.
\textit{arXiv:1407.3390}.

\bibitem[Clauset et al.(2009)]{clauset2009}
Clauset, A., Shalizi, C. R., \& Newman, M. E. J. (2009).
Power-law distributions in empirical data.
\textit{SIAM Review}, 51(4), 661--703.

\bibitem[Goliath \& Gebbie(2026)]{goliath2026}
Goliath, E., \& Gebbie, T. (2026).
Metaorder modelling and identification from public data.
\textit{arXiv:2602.19590}.

\bibitem[Lillo \& Farmer(2004)]{lillo2004}
Lillo, F., \& Farmer, J. D. (2004).
The long memory of the efficient market.
\textit{Studies in Nonlinear Dynamics and Econometrics}, 8(3).

\bibitem[Lillo et al.(2005)]{lillo2005}
Lillo, F., Mike, S., \& Farmer, J. D. (2005).
Theory for long memory in supply and demand.
\textit{Physical Review E}, 71(6), 066122.

\bibitem[Maitrier et al.(2025)]{maitrier2025}
Maitrier, G., Loeper, G., \& Bouchaud, J.-P. (2025).
Generating realistic metaorders from public data.
\textit{arXiv:2503.18199}.

\bibitem[Mandelbrot(1997)]{mandelbrot1997}
Mandelbrot, B. B. (1997).
\textit{Fractals and Scaling in Finance}.
Springer, New York.

\bibitem[Sato \& Kanazawa(2023)]{sato2023}
Sato, Y., \& Kanazawa, K. (2023).
Inferring microscopic financial information from the long memory in
market-order flow: A quantitative test of the Lillo--Mike--Farmer model.
\textit{Physical Review Letters}, 131, 197401.

\bibitem[Yang \& Zhang(2000)]{yangzhang2000}
Yang, D., \& Zhang, Q. (2000).
Drift-independent volatility estimation based on high, low, open, and close
prices.
\textit{Journal of Business}, 73(3), 477--491.

\end{thebibliography}

\end{document}
